\begin{thebibliography}{99}

\bibitem{AAK}
M. Abouzaid, D. Auroux and L. Katzarkov,
\emph{Lagrangian fibrations on blowups of toric varieties and mirror symmetry for hypersurfaces},
Publications Mathématiques de l'IHÉS \textbf{123} (2016), 199--282.
\url{https://doi.org/10.1007/s10240-016-0081-9}.

\bibitem{BSV}
M. S. Borman, N. Sheridan and U. Varolgunes,
\emph{Quantum cohomology as a deformation of symplectic cohomology},
Journal of Fixed Point Theory and Applications \textbf{24} (2022), article 48.
Proposition~2.10 (commuting Hamiltonians and adapted primitive).
\url{https://doi.org/10.1007/s11784-022-00965-6};
\url{https://arxiv.org/abs/2108.08391v2}.

\bibitem{CGblowup}
P. Chakravarthy and Y. Groman,
\emph{Almost toric fibrations on symplectic blow ups},
arXiv:2510.00994v1 (2025).
Theorem~1.1, Definition~2.2, Lemma~3.7.
\url{https://arxiv.org/abs/2510.00994v1}.

\bibitem{CGK3}
P. Chakravarthy and Y. Groman,
\emph{Almost toric fibrations on K3 surfaces via degenerations},
Advances in Mathematics \textbf{502} (2026), Part~B, 111138.
Definition~1.2, Theorems~1.6, 1.8, Corollary~1.9;
Lemma~2.13 is cited in the revised arXiv version of 28 May 2026.
\url{https://doi.org/10.1016/j.aim.2026.111138};
\url{https://arxiv.org/abs/2502.04304v3}.

\bibitem{DMT}
J.-P. Dufour, P. Molino and A. Toulet,
\emph{Classification des systèmes intégrables en dimension 2 et invariants des modèles de Fomenko},
Comptes Rendus de l'Académie des Sciences, Paris, Série~I, Mathématique
\textbf{318} (1994), no.~10, 949--952.
An English translation is available at
\url{https://webpages.csus.edu/mjv/DufourMolinoToulet.pdf}.

\bibitem{DHJH}
S. Dwivedi, J. Herman, L. C. Jeffrey and T. van den Hurk,
\emph{Hamiltonian Group Actions and Equivariant Cohomology},
SpringerBriefs in Mathematics, Springer, Cham, 2019.
Theorem~3.2 and Corollary~3.4, pp.~17--20 (symplectic neighborhood);
Chapter~10 (Duistermaat--Heckman).
\url{https://doi.org/10.1007/978-3-030-27227-2}.

\bibitem{Evans}
J. Evans,
\emph{Lectures on Lagrangian Torus Fibrations},
Cambridge University Press, Cambridge, 2023.
\url{https://arxiv.org/abs/2110.08643v4}.

\bibitem{EvansNotes}
J. Evans,
\emph{Lagrangian torus fibrations (notes in progress)},
8 February 2019.
Lemmas~6.1.5, 6.2.1 and \S6.3.
The numbering here refers to this preliminary set of notes, not to \cite{Evans}.
\url{https://jde27.uk/misc/ltf.pdf}.

\bibitem{EM}
J. D. Evans and M. Mauri,
\emph{Constructing local models for Lagrangian torus fibrations},
Annales Henri Lebesgue \textbf{4} (2021), 537--570.
Lemmas~2.16, 2.22, 3.1; Remarks~2.7, 2.19, 2.23; \S3.1;
Assumption~6.6, Proposition~6.13 and Theorem~6.15.
\url{https://doi.org/10.5802/ahl.80};
\url{https://arxiv.org/abs/1905.09229v3}.

\bibitem{EvansUrzua}
J. D. Evans and G. Urz\'ua,
\emph{Antiflips, mutations, and unbounded symplectic embeddings of rational homology balls},
Ann. Inst. Fourier \textbf{71} (2021), no.~5, 1807--1843.
\url{https://doi.org/10.5802/aif.3429}.

\bibitem{BHS}
J. Brendel, J. Hauber and J. Schmitz,
\emph{Semi-Local Exotic Lagrangian Tori in Dimension Four},
Ann. Inst. Fourier, online first (2026), 48 pp.
\url{https://doi.org/10.5802/aif.3789};
\url{https://arxiv.org/abs/2403.00408}.

\bibitem{FBP}
J. Fernández de Bobadilla and T. Pełka,
\emph{Fibrations by Lagrangian tori for maximal Calabi--Yau degenerations and beyond},
Communications on Pure and Applied Mathematics (2026), e70060,
published online 8 July 2026,
\url{https://doi.org/10.1002/cpa.70060}.
Theorem~6.14, Proposition~5.6, Example~5.8 and Remark~5.9 are cited from the final preprint,
arXiv:2312.13248v4, 9 July 2026.
\url{https://arxiv.org/abs/2312.13248v4}.

\bibitem{Fleming}
A. Fleming,
\emph{Kähler packings of projective complex manifolds},
European Journal of Mathematics \textbf{7} (2021), 1382--1400.
Proposition~4.6 and Construction~4.7.
\url{https://doi.org/10.1007/s40879-021-00491-8};
\url{https://arxiv.org/abs/1905.03140}.

\bibitem{Forstneric}
F. Forstnerič,
\emph{Euclidean domains in complex manifolds},
Journal of Mathematical Analysis and Applications \textbf{506} (2022), no.~1, 125660.
The Introduction recalls the holomorphic tubular-neighborhood theorem of Docquier--Grauert.
\url{https://doi.org/10.1016/j.jmaa.2021.125660};
\url{https://arxiv.org/abs/2104.03051v3}.

\bibitem{GollaMartelli}
M. Golla and B. Martelli,
\emph{Pair of pants decomposition of 4-manifolds},
Algebraic \& Geometric Topology \textbf{17} (2017), no.~3, 1407--1444.
\url{https://doi.org/10.2140/agt.2017.17.1407};
\url{https://arxiv.org/abs/1503.05839}.

\bibitem{GV}
Y. Groman and U. Varolgunes,
\emph{Locality of relative symplectic cohomology for complete embeddings},
Compositio Mathematica \textbf{159} (2023), no.~12, 2551--2637.
\S7.3, Lemma~7.19 and Proposition~7.21.
\url{https://doi.org/10.1112/S0010437X23007492};
\url{https://arxiv.org/abs/2110.08891v4}.

\bibitem{GrossSiebert}
M. Gross and B. Siebert,
\emph{From real affine geometry to complex geometry},
Annals of Mathematics (2) \textbf{174} (2011), no.~3, 1301--1428.
\url{https://doi.org/10.4007/annals.2011.174.3.1}.

\bibitem{Hatcher}
A. Hatcher,
\emph{Algebraic Topology},
Cambridge University Press, Cambridge, 2002.
\S\S0.2, 2.2 and 4.2 (CW complexes, cellular homology and homotopy equivalences).
\url{https://pi.math.cornell.edu/~hatcher/AT/ATpage.html}.

\bibitem{Hartshorne}
R. Hartshorne,
\emph{Algebraic Geometry},
Graduate Texts in Mathematics \textbf{52}, Springer, New York, 1977.
Chapter~V, Theorem~1.10 (the Nakai--Moishezon criterion for surfaces).
\url{https://doi.org/10.1007/978-1-4757-3849-0}.

\bibitem{Hattori}
A. Hattori,
\emph{Topology of $\mathbb C^n$ minus a finite number of affine hyperplanes in general position},
Journal of the Faculty of Science, University of Tokyo, Section~IA, Mathematics
\textbf{22} (1975), no.~2, 205--219.
\url{https://repository.dl.itc.u-tokyo.ac.jp/records/39761}.

\bibitem{IKM}
A. Izosimov, B. Khesin and M. Mousavi,
\emph{Coadjoint orbits of symplectic diffeomorphisms of surfaces and ideal hydrodynamics},
Annales de l'Institut Fourier \textbf{66} (2016), no.~6, 2385--2433.
Theorems~3.11 and 3.17 (measured Reeb graphs and weighted Reeb graphs).
\url{https://doi.org/10.5802/aif.3066};
\url{https://arxiv.org/abs/1504.05629}.

\bibitem{LeungSymington}
N. C. Leung and M. Symington,
\emph{Almost toric symplectic four-manifolds},
Journal of Symplectic Geometry \textbf{8} (2010), no.~2, 143--187.
Theorem~2.6 and Table~1.
\url{https://doi.org/10.4310/JSG.2010.v8.n2.a2};
\url{https://arxiv.org/abs/math/0312165}.

\bibitem{LeungWan}
N. C. Leung and T. Y. H. Wan,
\emph{Calabi--Yau components in general type hypersurfaces},
Journal of Differential Geometry \textbf{83} (2009), no.~1, 43--74.
\S\S1 and 3.1.
\url{https://doi.org/10.4310/jdg/1253804351};
\url{https://arxiv.org/abs/0807.1784}.

\bibitem{Maksymenko}
S. Maksymenko,
\emph{Path-components of Morse mappings spaces of surfaces},
Commentarii Mathematici Helvetici \textbf{80} (2005), no.~3, 655--690.
Main Theorem, Lemma~4.1 and \S6.3.
\url{https://doi.org/10.4171/CMH/30};
\url{https://arxiv.org/abs/math/9910085}.

\bibitem{MatessiPants}
D. Matessi,
\emph{Lagrangian pairs of pants},
International Mathematics Research Notices \textbf{2021} (2021), no.~15, 11306--11356.
\url{https://doi.org/10.1093/imrn/rnz126};
\url{https://arxiv.org/abs/1802.02993}.

\bibitem{Matessi}
D. Matessi,
\emph{Lagrangian submanifolds from tropical hypersurfaces},
International Journal of Mathematics \textbf{32} (2021), no.~7, 2150046, 63~pp.
\url{https://arxiv.org/abs/1804.01469}.

\bibitem{IvanSmithTorus}
I. Smith,
\emph{Torus fibrations on symplectic four-manifolds},
Turkish Journal of Mathematics \textbf{25} (2001), no.~1, 69--96.
\url{https://journals.tubitak.gov.tr/math/vol25/iss1/4/}.
分類の仮定は connected homologically essential Lagrangian tori と
tame singularities である（原論文のabstractと定理の射程を区別する）．

\bibitem{Mikhalkin}
G. Mikhalkin,
\emph{Decomposition into pairs-of-pants for complex algebraic hypersurfaces},
Topology \textbf{43} (2004), no.~5, 1035--1065.
Theorems~1 and~1$'$; \S4.2.
\url{https://doi.org/10.1016/j.top.2003.11.006};
\url{https://arxiv.org/abs/math/0205011v3}.

\bibitem{PapadimaSuciu}
S. Papadima and A. I. Suciu,
\emph{Higher homotopy groups of complements of complex hyperplane arrangements},
Advances in Mathematics \textbf{165} (2002), 71--100.
\S1.1 (Hattori's theorem for generic arrangements).
\url{https://arxiv.org/abs/math/0002251}.

\bibitem{Pascaleff}
J. Pascaleff,
\emph{On the symplectic cohomology of log Calabi--Yau surfaces},
Geometry \& Topology \textbf{23} (2019), no.~6, 2701--2792.
\S4.2.
\url{https://doi.org/10.2140/gt.2019.23.2701};
\url{https://arxiv.org/abs/1304.5298v3}.

\bibitem{Ruan}
W.-D. Ruan,
\emph{Lagrangian torus fibration of quintic hypersurfaces. I. Fermat quintic case},
in \emph{Winter School on Mirror Symmetry, Vector Bundles and Lagrangian Submanifolds}
(Cambridge, MA, 1999), AMS/IP Studies in Advanced Mathematics \textbf{23},
American Mathematical Society, Providence, RI, 2001, 297--332.
\url{https://arxiv.org/abs/math/9904012}.

\bibitem{RSTZ}
H. Ruddat, N. Sibilla, D. Treumann and E. Zaslow,
\emph{Skeleta of affine hypersurfaces},
Geometry \& Topology \textbf{18} (2014), no.~3, 1343--1395.
\url{https://doi.org/10.2140/gt.2014.18.1343};
\url{https://arxiv.org/abs/1211.5263}.

\bibitem{StacksAmpleness}
The Stacks Project Authors,
\emph{The Stacks Project}, 2026.
Lemma~30.17.2, Tag~0B5V (ampleness under finite surjective pullback);
Lemma~37.50.3, Tag~0D2S (relative ampleness near an ample fiber);
Lemma~71.15.3, Tag~0D3A (the algebraic-space version).
\url{https://stacks.math.columbia.edu/tag/0B5V};
\url{https://stacks.math.columbia.edu/tag/0D2S};
\url{https://stacks.math.columbia.edu/tag/0D3A}.

\bibitem{VuNgoc}
S. Vũ Ngọc,
\emph{On semi-global invariants for focus--focus singularities},
Topology \textbf{42} (2003), no.~2, 365--380.
Theorem~2.1 and \S3.
\url{https://doi.org/10.1016/S0040-9383(01)00026-X};
\url{https://arxiv.org/abs/math/0208245}.

\bibitem{Zhang}
Y. Zhang,
\emph{Pair-of-Pants Decompositions of 4-Manifolds Diffeomorphic to General Type Hypersurfaces},
Journal of Mathematical Study \textbf{58} (2025), no.~2, 211--252.
Theorem~1.1; the quintic example in \S1.
\url{https://doi.org/10.4208/jms.v58n2.25.06};
\url{https://arxiv.org/abs/2102.10037}.

\bibitem{Zhou}
P. Zhou,
\emph{Lagrangian skeleta of hypersurfaces in $(\mathbb C^*)^n$},
Selecta Mathematica (N.S.) \textbf{26} (2020), no.~2, article~26, 33~pp.
\url{https://arxiv.org/abs/1803.00320}.

% Bibliography for the SNC theorem and subsequent applications.

\bibitem{DGMS}
P. Deligne, P. Griffiths, J. Morgan and D. Sullivan,
\emph{Real homotopy theory of K\"ahler manifolds},
Inventiones Mathematicae \textbf{29} (1975), no.~3, 245--274.
Lemma~(5.11), pp.~266--267, and Remark~(5.14).
\url{https://doi.org/10.1007/BF01389853};
\url{https://www.math.stonybrook.edu/~dennis/publications/PDF/DS-pub-0024.pdf}.

\bibitem{TMZ}
M. Farajzadeh Tehrani, M. McLean and A. Zinger,
\emph{Normal crossings singularities for symplectic topology},
Advances in Mathematics \textbf{339} (2018), 672--748.
Definitions~2.8, 2.9, 2.11, 2.12 and Theorem~2.13.
\url{https://doi.org/10.1016/j.aim.2018.09.035};
\url{https://arxiv.org/abs/1410.0609v3}.

\bibitem{FeichtnerSturmfels}
Eva Maria Feichtner and Bernd Sturmfels,
\emph{Matroid polytopes, nested sets and Bergman fans},
Portugaliae Mathematica (N.S.) \textbf{62} (2005), no.~4, 437--468.
\href{https://arxiv.org/abs/math/0411260}{arXiv:math/0411260}.

\bibitem{ArdilaKlivans}
Federico Ardila and Caroline J. Klivans,
\emph{The Bergman complex of a matroid and phylogenetic trees},
Journal of Combinatorial Theory, Series B \textbf{96} (2006), no.~1, 38--49.
\href{https://doi.org/10.1016/j.jctb.2005.06.004}{doi:10.1016/j.jctb.2005.06.004}.
\href{https://arxiv.org/abs/math/0311370}{arXiv:math/0311370}.

\bibitem{Shaw}
Kristin M. Shaw,
\emph{A tropical intersection product in matroidal fans},
SIAM Journal on Discrete Mathematics \textbf{27} (2013), no.~1, 459--491.
\href{https://doi.org/10.1137/110850141}{doi:10.1137/110850141}.
\href{https://arxiv.org/abs/1010.3967v2}{arXiv:1010.3967v2}.

\bibitem{KapranovChow}
Mikhail M. Kapranov,
\emph{Chow quotients of Grassmannians. I},
in \emph{I.~M.~Gelfand Seminar},
Advances in Soviet Mathematics, vol.~16, part~2,
American Mathematical Society, Providence, RI, 1993, pp.~29--110.
\href{https://arxiv.org/abs/alg-geom/9210002}{arXiv:alg-geom/9210002}.

\bibitem{Aksnes}
E. Aksnes,
\emph{Tropicalization of curve arrangement complements and arroids},
Mathematica Scandinavica \textbf{131} (2025), no.~2.
\url{https://doi.org/10.7146/math.scand.a-156674};
\url{https://arxiv.org/abs/2404.14380v3}.

\bibitem{GKM}
A. Gathmann, M. Kerber and H. Markwig,
\emph{Tropical fans and the moduli spaces of tropical curves},
Compositio Mathematica \textbf{145} (2009), no.~1, 173--195.
\url{https://doi.org/10.1112/S0010437X08003837};
\url{https://arxiv.org/abs/0708.2268v2}.

\bibitem{MichalakReeb}
\L{}. P. Michalak,
\emph{Realization of a graph as the Reeb graph of a Morse function on a manifold},
Topological Methods in Nonlinear Analysis \textbf{52} (2018), no.~2, 749--762.
Definition~2.2 and Theorem~5.6.
\url{https://doi.org/10.12775/TMNA.2018.029};
\url{https://arxiv.org/abs/1805.06727v2}.

% Add this entry inside the main thebibliography environment.

\bibitem{ShawSurfaces}
Kristin Shaw,
\emph{Tropical surfaces},
arXiv:1506.07407v1 (2015), 43 pp.
Definition~3.1, pp.~15--16 (ambient integer affine transitions);
\S3.2 (boundary), \S4.2 (tropical summation).
\url{https://arxiv.org/abs/1506.07407v1};
\url{https://arxiv.org/pdf/1506.07407}.

\bibitem{JellRauShaw}
Philipp Jell, Johannes Rau, and Kristin Shaw,
\emph{Lefschetz $(1,1)$-theorem in tropical geometry},
\'Epijournal de G\'eom\'etrie Alg\'ebrique \textbf{2} (2018), Article~11, 27 pp.
Definitions~2.1, 3.1 and 5.1 (rational polyhedral spaces, the affine sheaf, and matroidal charts).
\url{https://doi.org/10.46298/epiga.2018.volume2.4126};
\url{https://epiga.episciences.org/4981/pdf};
\url{https://arxiv.org/abs/1711.07900}.

\bibitem{LRSChern}
Luc\'ia L\'opez de Medrano, Felipe Rinc\'on, and Kris Shaw,
\emph{Chern Classes of Tropical Manifolds},
arXiv:2309.00229v2 (2023), 30 pp.
Definitions~2.2--2.4, pp.~4--5 (extended integer affine maps, tropical manifold atlas, and boundary).
\url{https://arxiv.org/abs/2309.00229v2};
\url{https://arxiv.org/pdf/2309.00229}.

\bibitem{CartwrightComplexes}
Dustin Cartwright,
\emph{Tropical complexes},
Manuscripta Mathematica \textbf{163} (2020), 227--261.
Proposition~2.1; Lemma~2.2; Definitions~2.5, 2.8;
Proposition~2.10; Construction~3.2 and Remark~3.4.
The linked arXiv version is v4 (2019).
\url{https://doi.org/10.1007/s00229-019-01148-3};
\url{https://arxiv.org/abs/1308.3813v4};
\url{https://arxiv.org/pdf/1308.3813}.

% Existing EM entry can additionally mention Assumption 6.6,
% Proposition 6.13 and Theorem 6.15.  Their log Calabi--Yau hypotheses
% must remain explicit; they do not prove the present branched atlas.
% Existing VuNgoc entry can additionally mention Remark 3.1,
% which displays the focus--focus monodromy matrix.

\bibitem{ShawWerner}
K. Shaw and A. Werner,
\emph{On the birational geometry of matroids},
Combinatorial Theory \textbf{3} (2023), no.~2, Paper No.~17.
Theorem~7.2 (parallel connections and product Bergman fans).
\url{https://doi.org/10.5070/C63261996};
\url{https://arxiv.org/abs/2207.13639v1}.

\bibitem{AllermannRau}
L. Allermann and J. Rau,
\emph{First steps in tropical intersection theory},
Mathematische Zeitschrift \textbf{264} (2010), no.~3, 633--670.
Definition~3.4, Remark~3.5, Proposition~3.7;
Definition~6.1 and Construction~6.4.
\url{https://doi.org/10.1007/s00209-009-0483-1};
\url{https://arxiv.org/abs/0709.3705v3}.

\bibitem{SymingtonFour}
M. Symington,
\emph{Four dimensions from two in symplectic topology},
arXiv:math/0210033v1 (2002).
Proposition~2.15, \S4.3, equation~(4.10), and Proposition~4.14.
\url{https://arxiv.org/abs/math/0210033v1}.

\bibitem{QuinticArithmetic}
\emph{五次曲面の整数アフィン接着と高分母支持グラフ}，
未発表草稿，第25稿，2026年10月28日．
本稿の companion 論文．


\bibitem{RuddatAnnouncement2026}
H. Ruddat,
\emph{Lagrangian torus fibrations of Calabi--Yau hypersurfaces},
Institut Mittag--Leffler, seminar abstract, 30 June 2026.
Joint work with C. Y. Mak, D. Matessi and I. Zharkov, as stated in the abstract.
Announcement; accessed 28 October 2026.
\url{https://www.mittag-leffler.se/seminar/helge-ruddat-tba/}.

\bibitem{ZharkovSoftCells2026}
I. Zharkov,
\emph{Phase Tropical hypersurfaces and their cousins},
ICMS--Sofia Geometry Seminar, 15 April 2026; announcement posted 8 April 2026.
Seminar abstract; accessed 28 October 2026.
\href{https://icms.bg/phase-tropical-hypersurfaces-and-their-cousins-geometry-seminar-talk-by-ilia-zharkov/}{ICMS--Sofia: seminar announcement}.


\bibitem{VenugopalanWoodward2026}
S. Venugopalan and C. T. Woodward,
\emph{Tropical disk potential for almost toric manifolds},
arXiv:2604.03161v1, 3 April 2026.
Definitions~1.1--1.2, Theorem~1.3, and the discussion of the nonmonotone potential in \S1.
\url{https://arxiv.org/abs/2604.03161v1}.

\bibitem{ZhangNegativeKE}
Y. Zhang,
\emph{Collapsing of negative Kähler--Einstein metrics},
Mathematical Research Letters \textbf{22} (2015), no.~6, 1843--1869.
\url{https://doi.org/10.4310/MRL.2015.v22.n6.a16}.
Theorems~1.1, 2.4, 2.6 in \url{https://arxiv.org/abs/1505.04728v2}.

\bibitem{SongNegativeKE}
J. Song,
\emph{Degeneration of Kähler--Einstein manifolds of negative scalar curvature},
arXiv:1706.01518v1 (2017), Theorem~1.2.
\url{https://arxiv.org/abs/1706.01518}.

\bibitem{MandelNegativeKE}
H. Mandel,
\emph{Degenerations of negative Kähler--Einstein surfaces},
Journal of the London Mathematical Society \textbf{109} (2024), no.~1, e12818.
\url{https://doi.org/10.1112/jlms.12818}.

\bibitem{KECubicalCurrent}
K. Ueda,
\emph{多重スケール Kähler--Einstein 幾何 V：超平面和への退化，立方体極限，および直径},
現行未公表原稿．定理 \texttt{thm:degree-d-limit}，
仮定 \texttt{hyp:CPd}, \texttt{hyp:PBEd}．
\url{https://github.com/uedakazushi/kaehler-einstein-degeneration/blob/main/papers/cubical-gh-limit.tex}.

\bibitem{RuanHMinimal}
W.-D. Ruan,
\emph{H-minimal Lagrangian fibrations in Kähler manifolds and minimal
Lagrangian vanishing tori in Kähler--Einstein manifolds},
arXiv:math/0309177v2 (2003).
Proposition~3.1, Lemma~3.3, and Theorems~3.1, 4.2.
\url{https://arxiv.org/abs/math/0309177v2}.

\bibitem{KELocalCurrent}
K. Ueda,
\emph{多重スケール Kähler--Einstein 幾何 IV：超平面和への退化と局所模型},
現行未公表原稿．
\texttt{coarse:hyp:TC5}，
\texttt{coarse:prop:density-counterexample}，
\texttt{coarse:lem:actual-buffer-comparison}．
\url{https://github.com/uedakazushi/kaehler-einstein-degeneration/blob/main/papers/kaehler-einstein-degeneration.tex}.

\end{thebibliography}
